欧拉路恰好经过每条边一次，允许重复经过顶点；起点与终点相同的欧拉路称为欧拉回路。以下条件均为充要条件，允许重边和自环。

\textbf{无向图：}所有非零度顶点位于同一个连通块，且满足下列一种情况：
\begin{itemize}
    \item 所有顶点度数均为偶数：存在欧拉回路，可从任一非零度顶点出发。
    \item 恰有两个奇度数顶点：存在非闭合欧拉路，起点、终点分别为这两个顶点，顺序任意。
\end{itemize}
无向自环对度数的贡献为 $2$。

\textbf{有向图：}忽略边的方向后，所有入度与出度之和非零的顶点位于同一个连通块，且满足下列一种情况：
\begin{itemize}
    \item 每个顶点的入度等于出度：存在欧拉回路，可从任一出度非零的顶点出发。
    \item 恰有一个顶点出度比入度大 $1$，作为起点；恰有一个顶点入度比出度大 $1$，作为终点；其余顶点入度等于出度：存在非闭合欧拉路。
\end{itemize}
有向自环使入度、出度各增加 $1$。

孤立点不影响上述条件。没有边时，任一顶点都可作为长度为 $0$ 的欧拉回路的起点。
